% !TeX root = ../../main.tex
\section{\bilingual{Conclusions}{结论}}
\label{section:conclusion-and-prospects}

\bilingual{
    This thesis studies the asymptotic evolutionary behavior of $N$-branching random walks with general branching laws under the framework of universality class with finite variance ($\alpha=2$). By introducing a second moment constraint for the cross terms of the branching law (\Cref{assumption:log-expt-cross-sum-exp-neg-sum}) and a moment constraint for the leftmost particle (\Cref{assumption:moment-exp-1}), we obtained the following main conclusions:
}{
    本文在有限方差 ($\alpha=2$) 的普适性类框架下, 研究了具有一般分支律的 $N$-分支随机游走的渐近演化行为. 通过引入针对分支律交叉项的二阶矩约束 (\Cref{assumption:log-expt-cross-sum-exp-neg-sum}) 及最左侧粒子的矩约束 (\Cref{assumption:moment-exp-1}) , 本文得到了以下主要结论:
}

\begin{enumerate}
    \item
    \bilingual{
        \textbf{Universality Extension of the Trajectory Limit Theorem}:

        This paper extends the $L^2$ limit theorem for the leftmost particle position, originally established by A\"id\'ekon and Hu~\mycite{aïdékon2025averageweightpercolationcompletegraph} for a specific Poisson point process branching law, to branching laws satisfying general moment conditions. Under the scaling where time is of order $t\log^3 N$ and space is of order $\log N$, the leftmost position of the $N$-BRW converges in $L^2$ to the line $-a+\frac{\pi^2}2\sigma^2t$, pointwise at each time slice $t$, as $N\to\infty$. This result quantitatively reveals how the initial particle count $\lfloor N^a\rfloor$ governs the frontier trajectory and confirms that the linear trajectory is an intrinsic feature of finite-variance $N$-BRW systems.
    }{
        \textbf{轨迹极限定理的普适性推广}:%
        本文将 A\"id\'ekon 与 Hu \mycite{aïdékon2025averageweightpercolationcompletegraph} 针对特定 Poisson 分支律建立的 $L^2$ 极限定理推广至一般矩条件情形. 当时间尺度取 $t\log^3 N$、空间尺度取 $\log N$ 时, $N$-BRW 的最左侧位置在每个时间切片 $t$ 处依 $L^2$ 收敛到直线 $-a+\frac{\pi^2}2\sigma^2t$, 从而定量刻画了初始粒子数 $\lfloor N^a\rfloor$ 对前沿演化轨迹的影响, 并表明该线性轨迹是有限方差 $N$-BRW 系统的本质特征.
    }

    \item
    \bilingual{
        \textbf{Application and Validation of the Mogul'ski\u i Small Deviation Estimation Framework}:

        We follow the small-deviation method of A\"id\'ekon and Hu~\mycite{aïdékon2025averageweightpercolationcompletegraph}. Lemmas \Cref{lemma:4.1,lemma:4.2} bound walks in horizontal and tilted tubes from above, and \Cref{lemma:4.3} gives the lower bound. Combined with exponential change of measure and \nameref{theorem:mogulskii}, these estimates support second-moment arguments for general branching laws.
    }{
        \textbf{\nameref{theorem:mogulskii}小偏差估计框架的应用与验证}:
        本文采用 A\"id\'ekon 与 Hu \mycite{aïdékon2025averageweightpercolationcompletegraph} 的小偏差方法. 引理 \Cref{lemma:4.1,lemma:4.2} 给出水平与倾斜管道的上界, 引理 \Cref{lemma:4.3} 给出下界. 结合指数测度变换和\nameref{theorem:mogulskii}, 这些估计支持一般分支律下的二阶矩分析.
    }

    \item
    \bilingual{
        \textbf{Technical Improvement in the $a=0$ Case Proof}:

        For the limit theorem starting from a single particle ($a=0$), this paper refines the proof of A\"id\'ekon and Hu~\mycite{aïdékon2025averageweightpercolationcompletegraph}. Using the supercritical assumption, a general branching law is coupled to an at-most-binary branching process via displacement truncation $M$, and a {\kaishu restart mechanism} combined with submartingale maximal inequalities controls position deviations, thereby avoiding the extinction risk induced by truncation. Under this construction the population reaches size $O(N^\alpha)$ in $O(\log N)$ time, ensuring stability of the limit trajectory at the leading scale.
    }{
        \textbf{$a=0$ 情形证明的技术改进}:%
        对于从单个粒子出发 ($a=0$) 的极限定理, 本文改进了 A\"id\'ekon 与 Hu \mycite{aïdékon2025averageweightpercolationcompletegraph} 的证明. 利用超临界假设, 通过位移截断 $M$ 将一般分支律耦合成至多二分叉过程, 并引入{\kaishu 重启机制}配合下鞅最大值不等式控制位置偏移, 从而规避了一般分支律截断后的灭绝风险. 该构造下种群规模达到 $O(N^\alpha)$ 仅需 $O(\log N)$ 时间, 保证了极限轨迹在主尺度下的稳定性.
    }

    \item \bilingual{
        \textbf{Comparative Evaluation with Mallein's Conclusion}:

        Compared with Mallein~\mycite{Mallein_2017} on the speed $v_N$, this paper obtains stronger $L^2$ convergence at the cost of narrower distributional universality (requiring exponential moments). For $\alpha\in(0,2]$, Mallein established the almost-sure limit of $v_N$ under second moment conditions; for $\alpha=2$, the present paper, via finer second moment analysis, confirms the $(\log N)^{-2}$ scaling of the velocity correction with an explicit constant and characterizes the full dynamic trajectory of the front during equilibration. The one-sided endpoint estimate proves the lower-speed bound in the first-moment regime; the matching upper bound in \Cref{theorem:speed-L1} remains open.
    }{
        \textbf{与 Mallein 结论的对比}:%
        与 Mallein \mycite{Mallein_2017} 关于速度 $v_N$ 的结果相比, 本文以更窄的分布普适性为代价 (需指数矩控制方差), 换取了更强的 $L^2$ 收敛. ~Mallein 在二阶矩条件下对 $\alpha\in(0,2]$ 确立了 $v_N$ 的几乎处处极限; 本文针对 $\alpha=2$, 通过更精细的二阶矩分析验证了速度修正的 $(\log N)^{-2}$ 标度律并给出显式常数, 同时完整刻画了前沿在平衡建立过程中的动态轨迹. 一阶矩条件下的一侧端点估计已给出相应的速率下界; \Cref{theorem:speed-L1} 中匹配的上界仍待证明.
    }
\end{enumerate}
